% =============================================================================
% ch02_related.tex  —  Chapter 2: Background and Related Work
% K-Bound long-form monograph.
% NO \documentclass or \usepackage. Begins with \chapter.
% =============================================================================

\chapter{Background and Related Work}\label{ch:related}

Five lines of work surround K-Bound, and each comes within reach of its
question without asking it. The recurring distinction --- the one the theory
formalizes --- is between \emph{improving} an unlabeled objective and
\emph{certifying the sign} of its benefit before acting. We walk the five lines
in that order, and the gap widens as we go: TTA mechanisms adapt without asking
whether they should; the failure literature shows that they often should not;
guarded adaptation reacts to harm once it has begun; label-free estimation and
identifiability theory ask the right question but answer a harder one or stop at
asymptotics; selective prediction and conformal inference supply the abstention
and interval machinery but not the decision. By the final section every piece K-Bound
needs has appeared in isolation; what no prior work assembles is the single
pre-adaptation adapt/freeze/abstain certificate with a proven unknowable region
and a false-adapt bound. We close with a positioning table and a precise
statement of that gap.

\section{Test-Time Adaptation: Mechanisms and Methods}\label{sec:rel-tta}

Test-time adaptation originated as a practical answer to the distribution-shift
problem: rather than retraining, use the unlabeled target stream to adjust the
deployed model.
The field has developed along several distinct axes, each representing a
different hypothesis about what makes TTA fragile and how to make it robust.

\paragraph{Entropy minimization.}
\textbf{Tent} \cite{wang2021tent} is the canonical entropy-minimization method.
At each test batch, Tent minimizes the Shannon entropy of the model's softmax
predictions over only the batch-normalization affine parameters, keeping the
rest of the network frozen.
The intuition is that predictions on target data should be confident
(low entropy), and adjusting normalization statistics accomplishes this with
minimal risk of catastrophic forgetting.
On i.i.d.\ corruption benchmarks such as CIFAR-10-C
\cite{hendrycks2019benchmarking} and ImageNet-C, Tent yields substantial
accuracy improvements over the frozen baseline.

The fragility of pure entropy minimization, however, is well documented.
Tent applied continually to non-stationary or class-imbalanced streams
accumulates normalization-parameter damage; once the predictions collapse to
a single class, entropy continues to decrease while accuracy plummets.
Our stress-grid experiments (Chapter~\ref{ch:exp2}) reproduce this collapse
in the harm-dominated online CIFAR-10 regime, where always-adapt with Tent
falls to $0.339$ accuracy while a frozen model holds $0.440$.

\paragraph{Batch-normalization statistics adaptation.}
\textbf{Schneider et al.}~\cite{schneider2020bn} showed that a large fraction
of the performance gap on corruption benchmarks can be closed simply by
re-estimating batch-normalization running statistics from a mixture of source
and target data, without any gradient update.
This observation motivated a family of statistics-only approaches that avoid
the optimization risk of entropy minimization entirely, at the cost of needing
a sufficient target batch size to estimate stable statistics.
\textbf{NOTE}~\cite{gong2022note} extends this direction to non-i.i.d.\ streams
using class-conditional normalization, addressing the label-shift fragility of
global statistics.

\paragraph{Filtering and Fisher regularization.}
\textbf{EATA}~\cite{niu2022eata} addresses Tent's catastrophic-forgetting
tendency with two additions: an entropy filter that rejects unreliable
(high-entropy) samples from the update, and a Fisher-information-based
regularization penalty that penalizes updates proportional to their importance
for source-domain performance.
The entropy filter prevents the worst collapse triggers from entering the
update gradient; the Fisher penalty ensures the normalization parameters do
not drift too far from their source-optimal values.
EATA retains substantial improvement over the frozen baseline on standard
benchmarks while partially protecting against the forgetting that pure Tent
suffers.

\paragraph{Sharpness-aware updates and resets.}
\textbf{SAR}~\cite{niu2023sar} adds sharpness-aware entropy minimization
combined with an online reset mechanism.
Sharpness-aware optimization minimizes an upper bound on the worst-case
perturbation of the loss landscape, producing parameter updates that land in
flat regions less susceptible to subsequent drift.
The reset mechanism monitors a model-divergence criterion and resets the
normalization parameters to their source values if divergence exceeds a
threshold, aborting incipient collapse.
SAR is notably more robust than Tent or EATA on the continual and
non-stationary regimes that provoke collapse; our experiments confirm this,
with SAR's harmful base rate ($16\%$ on the stress grid) substantially lower
than Tent's ($34\%$) and EATA's ($27\%$).
This robustness is also the reason K-Bound ties rather than beats SAR:
when the candidate already self-protects, the certificate's additional value
is limited.

\paragraph{Source-free and continual adaptation.}
\textbf{SHOT} trains a source-free adaptation by information maximization and
self-supervised pseudo-labeling, enabling adaptation when the source data is
unavailable at deployment.
\textbf{CoTTA}~\cite{wang2022cotta} addresses the continual adaptation setting
by stochastic restoring: with some probability, parameters are reset to their
source-pre-trained values, preventing irreversible forgetting across a long
non-stationary stream.
\textbf{TTT}~\cite{sun2020ttt} augments the source training with a
self-supervised auxiliary task and continues training that task at test time,
adapting the shared feature extractor.
\textbf{MEMO}~\cite{zhang2022memo} takes a different route: for each test
point it generates multiple augmented views and minimizes the marginal entropy
across augmentations, adapting per-sample rather than per-batch.

\paragraph{Parameter-free approaches.}
\textbf{LAME}~\cite{boudiaf2022lame} is at the opposite extreme: it requires
no parameter update at all, instead post-processing the model's predictions
using a Laplacian regularization that encourages smoothness across the target
batch.
Its robustness to collapse is guaranteed by construction, since there are no
parameters to corrupt, but it cannot recover the same accuracy gains as
gradient-based methods when the distribution shift is large.

\paragraph{What all TTA methods share.}
Every method in this section, whether gradient-based or gradient-free, whether
it uses resets or regularization or filtering, makes the same architectural
commitment: \emph{it adapts}.
The question of whether adaptation is the right action on a given deployment
batch is not posed; the decision is implicit in the method's deployment.
K-Bound is mechanism-agnostic with respect to this entire class of methods.
It wraps any such candidate --- we evaluate Tent, EATA, and SAR --- and
makes the adapt/freeze/abstain decision at the certificate level, before
the adaptation mechanism executes.

\section{TTA Failure, Collapse, and Evaluation}\label{sec:rel-failure}

If the previous section's methods all commit to adapting, the next question is
how often that commitment backfires. A parallel line of work answers it: by
cataloguing the conditions under which TTA fails destructively and showing that
those conditions are not exotic but the realistic ones. This literature supplies
the empirical motivation for K-Bound's central premise --- that harmful adaptation
is common enough to be worth certifying against --- and the stress-test design
that our experiments inherit.

\paragraph{Pitfalls and failure modes.}
\textbf{Zhao et al.}~\cite{zhao2023pitfalls} systematically catalogued the
pitfalls of test-time adaptation evaluation: sensitivity to the batch-size
regime, stream composition, update budget, and the distinction between
i.i.d.\ and temporally correlated streams.
Their analysis showed that many published TTA results are obtained in regimes
where adaptation is effectively guaranteed to help (large i.i.d.\ batches of
a single corruption type), and that performance degrades substantially ---
sometimes below the frozen baseline --- in the more realistic regimes of mixed,
imbalanced, or small-batch streams.
This work directly motivates the design of the K-Bound pre-registered stress
grid: we cross severity, batch size, composition, and aggressiveness to
construct conditions that span the helpful and harmful regimes.

\paragraph{Robustness benchmarks.}
\textbf{RDumb}~\cite{press2023rdumb} demonstrated that a periodically
resetting stupid baseline (no adaptation, periodic reset to the source model)
outperforms many published TTA methods on long, non-stationary streams.
The implication is that sophisticated adaptation mechanisms may be
over-optimized for the benign conditions under which they are evaluated and
break down precisely when deployment conditions are most challenging.

\paragraph{Standardized benchmarks.}
\textbf{Hendrycks and Dietterich}~\cite{hendrycks2019benchmarking} introduced
the CIFAR-10-C, CIFAR-100-C, and ImageNet-C benchmarks, which apply 19
algorithmically defined corruptions at five severity levels to standard test
sets.
These benchmarks have become the standard evaluation platform for TTA methods;
they are also the platform for our per-corruption experiment
(Chapter~\ref{ch:exp2}).
However, as Zhao et al.\ and Press et al.\ note, per-corruption evaluation
is inherently helpful-dominated: each batch contains a single corruption type
at a single severity, giving the adaptation mechanism a maximally informative
signal.
Our stress grid goes beyond this by mixing compositions.

\paragraph{DomainBed.}
\textbf{Gulrajani and Lopez-Paz}~\cite{gulrajani2021domainbed} introduced the
DomainBed benchmark for domain generalization, emphasizing standardized
evaluation protocols and a thorough comparison of algorithms under matched
hyperparameter selection.
Their finding that many domain generalization methods do not consistently
outperform empirical risk minimization under rigorous evaluation is
methodologically analogous to our finding about TTA in the benign regime:
sophisticated methods add value primarily in the regimes where they are
designed to help.
\textbf{WILDS}~\cite{koh2021wilds} extends this to real-world distribution
shifts, including medical imaging, satellite imagery, and user-generated
content.
Extending K-Bound to WILDS benchmarks is an explicit item of future work
(Chapter~\ref{ch:discussion}).

\section{Guarded and Protected Adaptation}\label{sec:rel-guarded}

Given that adaptation can collapse, a natural response is to guard it. The
methods in this section --- distinct from both the mechanisms above and the
failure analyses just surveyed --- protect adaptation against specific failure
modes \emph{during} the update. They are the closest prior work to K-Bound in
spirit, which makes the one structural difference between them and us worth
stating precisely: a guard acts while adapting, whereas K-Bound decides before
adapting at all.

\paragraph{Entropy thresholding and filtering.}
EATA's entropy filter is the most direct example: high-entropy samples are
excluded from the update gradient on the hypothesis that they carry
corrupted or uninformative signal.
This protects against one failure direction (adapting on adversarial inputs)
but does not address the three-way decision between helpful, harmful, and
unknowable evidence, nor does it provide a formal bound on the false-adapt
rate.

\paragraph{Betting-style monitors.}
Sequential testing literature has developed monitors that halt or dampen a
running process when a specific failure mode is detected.
In the TTA setting, such a monitor could, for example, halt entropy
minimization when the estimated class distribution collapses below a threshold.
These approaches protect one direction --- typically against collapse --- but
they do not pose the three-way decision with an explicit unknowable region,
and they do not certify \emph{before} acting that adaptation will help.

\paragraph{NOTE's class-conditional normalization.}
\textbf{NOTE}~\cite{gong2022note} addresses the temporal correlation between
samples in a non-stationary stream by estimating class-conditional
normalization statistics.
This is a form of guarded adaptation in the sense that it conditions the
normalization update on the estimated class label of each sample, reducing
the sensitivity to class imbalance that causes Tent to collapse.
The guard is built into the mechanism, not into a separate decision layer.

\paragraph{The structural gap.}
All guarded-adaptation approaches share the architecture of operating
\emph{during} the update.
K-Bound operates \emph{before} the update: it makes the adapt/freeze/abstain
decision first, and only then --- if the decision is \textsc{adapt} ---
does it invoke the underlying mechanism.
This separation is not merely organizational; it is what makes the
false-adapt guarantee meaningful.
A certificate that fires after an adaptation has already occurred cannot
meaningfully bound the false-adapt rate; one that fires before can.

\section{Label-Free Accuracy Estimation}\label{sec:rel-estimation}

Guarding adaptation presumes one can tell when it is going wrong, which raises a
sharper question: can target performance be read off unlabeled data directly? A
third line of work does exactly this --- predicting accuracy from unlabeled
inputs --- but stops short of turning the prediction into an adapt/freeze
decision, and estimates an absolute risk where K-Bound needs only the sign of a
difference between two models.

\paragraph{ATC.}
\textbf{Garg et al.}~\cite{garg2022atc} introduced Average Threshold
Confidence (ATC), which thresholds source-calibrated softmax confidence to
estimate the fraction of correct predictions on a target set.
ATC achieves strong empirical accuracy on standard benchmarks and has become
a reference method for label-free performance estimation.
Its training requires a labeled source set to calibrate the confidence
threshold; given that calibration, it requires only unlabeled target data.

\paragraph{Agreement-on-the-Line.}
\textbf{Baek et al.}~\cite{baek2022agreement} exploited the empirical
observation that model agreement (the fraction of inputs on which two
independently trained models make the same prediction) correlates linearly
with accuracy on both in-distribution and out-of-distribution data.
The linear relationship can be calibrated from in-distribution data and
extrapolated to estimate OOD accuracy from OOD agreement, which is label-free
and observable.
The method is particularly valuable in the covariate-shift regime where the
linear correlation holds reliably, and less so in label-shift regimes where
it breaks down.

\paragraph{Predicting with confidence.}
\textbf{Guillory et al.}~\cite{guillory2021predicting} studied the problem
of estimating model accuracy on unseen distributions with confidence
intervals, developing methods based on source/target distribution divergence
and ensemble disagreement.

\paragraph{The distinction from K-Bound.}
All of these methods estimate $R_T(f)$ or $R_T(f_0)$ --- the target risk of
a single model --- or detect that TTA has failed \emph{post hoc}.
K-Bound's object of study is strictly different: the \emph{sign} of
$R_T(f_0) - R_T(f_a)$.
This is a relational quantity involving \emph{two} models, not an absolute
risk estimate.
Certifying the sign is a strictly weaker requirement than estimating either
absolute risk, so the knowable set is strictly larger --- but, as
Theorem~\ref{thm:imp} establishes, it is not all of evidence space.
Furthermore, K-Bound provides a \emph{pre-adaptation} certificate with a
controlled false-adapt rate: it fires before $f_a$ is deployed, not after.
Post-hoc detection that TTA has hurt (as in AETTA and some ensemble-based
monitors) is useful for retrospective analysis but cannot prevent a false
adaptation from occurring.

\section{Unsupervised Risk Estimation and Identifiability}\label{sec:rel-ident}

The estimators above succeed in some shift regimes and fail in others, and the
reason is not engineering but identifiability: whether the target risk is a
function of what unlabeled data can reveal at all. This is the line that gives
K-Bound both its hardness and its escape route. The impossibility results here
are the ancestors of our unknowable floor; the covariate-shift and label-shift
identifiability results are the ancestors of our positive regimes. K-Bound's move
within this line is to ask for the sign of a risk difference rather than a risk,
which is identifiable on a strictly larger set --- and then to make that
asymptotic fact a finite-sample certificate.

\paragraph{Steinhardt and Liang.}
Steinhardt and Liang showed that unsupervised risk estimation is possible only
under a conditional-independence structure between the features, the label,
and the observed predictions.
In particular, they required that the label can be recovered from the
prediction via a known conditional distribution, which is essentially an
assumption of calibration stability.
Under this assumption, the target risk is identifiable from unlabeled data;
without it, the problem is generally intractable.
Their work establishes the structural condition under which label-free risk
estimation is possible and motivates our weaker target: the sign rather
than the magnitude of a risk difference.

\paragraph{Ben-David et al.}
\textbf{Ben-David et al.}~\cite{bendavid2010theory} established the canonical
learning-theoretic framework for domain adaptation.
Their central bound expresses the target risk as the source risk plus an
$\mathcal{H}$-divergence between the source and target marginals plus a
residual term capturing the irreducible concept-shift component.
The $\mathcal{H}$-divergence is computable from unlabeled data; the residual
term is not.
This bound is tight in the sense that all three terms are necessary; the
identifiability result of Theorem~\ref{thm:pos} can be understood as showing
that under covariate shift the residual term vanishes and the $\mathcal{H}$-divergence
is controlled by the density ratio.

\paragraph{Mansour et al.}
\textbf{Mansour et al.}~\cite{mansour2009domain} studied domain adaptation
from a learning-bounds perspective, extending the Ben-David et al.\ framework
to mixture-of-source distributions and developing discrepancy-based bounds.
Their discrepancy measure generalizes the $\mathcal{H}$-divergence and admits
importance-weighting interpretations that connect to our covariate-shift
positive result.

\paragraph{Label-shift correction.}
\textbf{Lipton et al.}~\cite{lipton2018detecting} developed Black Box Shift
Estimation (BBSE), which corrects for label shift --- a change in the marginal
$P(Y)$ with stable class-conditional distributions $P(X \mid Y)$ --- using
an invertible confusion matrix estimated from labeled source data.
BBSE identifies the target label marginal from unlabeled target predictions,
enabling unbiased risk estimation under label shift.
This is one of the two identifiable special cases in the knowability hierarchy
(Table~\ref{tab:rel} in Chapter~\ref{ch:formulation}): label shift with an
invertible confusion matrix places us in the knowably helpful or knowably
harmful regime, not the unknowable one.

\paragraph{Dataset shift taxonomy.}
\textbf{Quino\~nero-Candela et al.}~\cite{quinonero2009dataset} provide the
standard taxonomy of dataset shift: covariate shift, label shift, concept
shift, and their combinations.
This taxonomy is the organizational backbone of the knowability hierarchy:
covariate shift is provably identifiable (Theorem~\ref{thm:pos}), label shift
is identifiable under a mild structural assumption (invertible confusion),
concept shift re-enters the unknowable regime (Theorem~\ref{thm:imp}), and
mixed/unknown shift requires certify-or-abstain.

\paragraph{Importance weighting.}
\textbf{Shimodaira}~\cite{shimodaira2000improving} introduced importance
weighting (weighted log-likelihood) as the canonical approach to covariate-shift
adaptation, showing that reweighting the source loss by the density ratio
$p_T(x)/p_S(x)$ yields consistent estimators of target risk.
\textbf{Sugiyama et al.}~\cite{sugiyama2007covariate} developed importance
weighted cross-validation, which uses the density ratio to select hyperparameters
in a way that is consistent under covariate shift.
These works are the direct antecedents of Theorem~\ref{thm:pos}: the
importance-weighted estimator $\widehat{R}_T(f) = n^{-1} \sum_i r(x_i)
\ell(f(x_i), y_i)$ is the operational mechanism behind the covariate-shift
positive result.

\paragraph{The K-Bound delta over this line.}
The identifiability literature establishes when risk is identifiable and when
it is not.
K-Bound's contribution within this line is twofold.
First, we target the \emph{sign} of a risk difference, which is identifiable
on a strictly larger set than the risk itself: the unknowable regime in
K-Bound is strictly smaller than the non-identifiable regime for risk
estimation.
Second, we instantiate the identifiability result as a \emph{finite-sample,
false-adapt-controlled certificate} (Theorem~\ref{thm:cert}), going from an
asymptotic identifiability statement to an operational, finite-sample
decision procedure.

\section{Selective Prediction and Abstention}\label{sec:rel-selective}

Identifiability theory tells us that some conditions are undecidable; acting on
that fact requires a principled way to decline. Abstention is not new to machine
learning --- it has been studied for decades --- but always at the granularity of
the individual prediction. K-Bound borrows the decision-theoretic logic of the
reject option and relocates it one level up, from the per-sample label to the
per-episode adaptation decision.

\paragraph{Chow's reject option.}
\textbf{Chow}~\cite{chow1970optimum} introduced the reject option for
classifiers: rather than committing to a prediction for every input, a
classifier may abstain when its confidence is below a threshold, trading
coverage for accuracy.
The optimality theory shows that the Bayes-optimal reject rule thresholds the
posterior probability, and the coverage-accuracy tradeoff is characterized by
the Bayes optimal.

\paragraph{Selective classification.}
\textbf{El-Yaniv and Wiener}~\cite{elyaniv2010selective} and
\textbf{Geifman and El-Yaniv}~\cite{geifman2017selective} developed the
modern theory of selective classification, providing finite-sample guarantees
on the accuracy of a classifier conditioned on its decision to predict.
SelectiveNet extends selective classification to deep neural networks with
an end-to-end trainable coverage mechanism.

\paragraph{Conformal prediction.}
\textbf{Vovk et al.}~\cite{vovk2005algorithmic} introduced conformal
prediction as a framework for constructing prediction sets with valid marginal
coverage.
\textbf{Angelopoulos and Bates}~\cite{angelopoulos2023conformal} provide a
recent tutorial treatment, including split-conformal prediction (which
uses a held-out calibration set to set the threshold) and risk-controlling
prediction sets.
Split-conformal prediction is the mechanism behind Theorem~\ref{thm:cert}:
the conformal radius $\varepsilon$ is constructed from a held-out set of
labeled source/validation benefits, and its coverage guarantee transfers to
the certificate.

\paragraph{K-Bound's abstention is over a decision, not a prediction.}
Selective prediction abstracts on individual \emph{predictions}: given a single
test input $x$, should the model commit to a class label $\hat{y}$ or abstain?
K-Bound's abstention is at a different level: given a deployment batch and an
observable evidence vector $Z$, should the system commit to the
\textsc{adapt} or \textsc{freeze} \emph{action}, or abstain from the
adaptation decision?
The object is not a label for a single point but a binary action for an entire
deployment episode.
This distinction matters because the cost of a wrong action is not the
per-sample misclassification rate but the regret relative to the oracle
policy --- an expectation over the entire deployed batch.

\section{Conformal Prediction and Anytime-Valid Inference}\label{sec:rel-conformal}

The reject option says \emph{when} to abstain; it does not say how to build the
calibrated interval whose width triggers the abstention. That interval is the
last component K-Bound draws from prior work. The conformal and anytime-valid
literature provides it --- finite-sample coverage for the batch certificate, and
valid-at-any-stopping-time coverage for its streaming extension --- with no claim
on the adapt/freeze/abstain framing itself.

\paragraph{Game-theoretic probability.}
\textbf{Ville}~\cite{ville1939etude} introduced the concept of e-processes
(nonneg.\ supermartingales) and established Ville's inequality, which gives
the anytime-valid analogue of Markov's inequality.
\textbf{Shafer and Vovk}~\cite{shafer2019game} developed the game-theoretic
foundations for probability and statistics, showing that betting-style
arguments yield valid inference without relying on the law of large numbers.

\paragraph{Testing by betting and e-values.}
\textbf{Shafer}~\cite{shafer2019game} and \textbf{Ramdas et al.}\
\cite{ramdas2023gametheoretic} developed the testing-by-betting framework:
a hypothesis is rejected when the cumulative product of e-values (factors
$1 + \lambda_t X_t$ where $X_t$ is a bounded observation) exceeds $1/\alpha$.
The product is a nonneg.\ supermartingale under the null, so its exceedance
is controlled at level $\alpha$ uniformly over all stopping times (Ville's
inequality), without any correction for repeated looks.

\paragraph{Anytime-valid confidence sequences.}
\textbf{Howard et al.}~\cite{howard2021timeuniform} constructed
time-uniform, nonparametric confidence sequences: intervals that are valid
at all sample sizes simultaneously.
\textbf{Waudby-Smith and Ramdas}~\cite{waudbysmith2024betting} developed
betting-based confidence intervals for bounded random variables, establishing
tight bounds via predictable bets.
\textbf{Gr\"unwald et al.}~\cite{grunwald2024safe} formalized safe testing,
which uses e-values to test composite hypotheses in a way that is valid
under optional stopping.

\paragraph{Connection to K-Bound.}
Theorem~\ref{thm:cert} is a fixed-$n$ certificate; the anytime-valid
extension (Appendix~\ref{app:repro}) uses a testing-by-betting e-process.
The anytime-valid version is strictly more powerful than any fixed-$n$
certificate in the sequential regime (it is valid at any stopping time with
no multiplicity correction), while the fixed-$n$ version is tighter at a
known sample size.
The empirical-Bernstein bound of \textbf{Maurer and Pontil}
\cite{maurer2009empirical} provides the radius discharge: for bounded
paired benefits in $[a,b]$, the empirical variance proxy yields a confidence
radius under independence/exchangeability.

\paragraph{Classical concentration.}
\textbf{Hoeffding}~\cite{hoeffding1963probability} established the classic
exponential tail bound for bounded independent random variables, and
\textbf{Le Cam}~\cite{lecam1973convergence} developed the two-point lower
bound technique (used in Theorem~\ref{thm:imp}) and the fundamental
convergence-of-estimates theory.
\textbf{Tsybakov}~\cite{tsybakov2009introduction} and
\textbf{Lehmann and Romano}~\cite{lehmann2005testing} are the standard
references for nonparametric estimation theory and testing, respectively.

\section{Benchmarks, Robustness, and Distribution Shift}\label{sec:rel-bench}

The five lines above frame the problem; this section names the datasets and
backbones on which we test it, grouped by the role each plays in the evaluation
of Chapters~\ref{ch:exp1} and~\ref{ch:exp2}.

\paragraph{Robustness benchmarks.}
\textbf{Hendrycks and Dietterich}~\cite{hendrycks2019benchmarking} introduced
the corruption benchmarks that anchor all of our CIFAR-10-C experiments.
\textbf{Recht et al.}~\cite{recht2019imagenet} showed that ImageNet classifiers
lose substantial accuracy on a re-collected ImageNet test set, motivating the
study of natural distribution shift.
\textbf{Taori et al.}~\cite{taori2020measuring} measured robustness to a broad
set of natural and synthetic distribution shifts, finding that effective
robustness (robustness beyond what is predicted by the in-distribution to
OOD correlation) is rare.
\textbf{Miller et al.}~\cite{miller2021accuracy} established the accuracy-on-the-line
phenomenon: across many models and distribution shifts, in-distribution and
OOD accuracy are linearly correlated.

\paragraph{Anomaly detection benchmarks.}
K-Bound's anomaly-routing experiments use the ADBench benchmark
\cite{han2022adbench}, which covers tabular, image, text, cyber, and fraud
anomaly detection across 57 datasets and multiple detector families.
\textbf{MVTec AD}~\cite{bergmann2019mvtec} and
\textbf{MVTec 3D-AD}~\cite{bergmann2022mvtec3d} provide industrial
visual-inspection benchmarks where distribution shift is induced by modality
failure; the anonymized multimodal instantiation (Chapter~\ref{ch:elara}) is evaluated
on MVTec 3D-AD.
\textbf{Real-IAD}~\cite{wang2024realiad} is a large-scale real industrial
anomaly detection dataset with 30 object categories; Real-IAD-D3 is the
anonymized multimodal reliability experiment.
\textbf{PyOD}~\cite{zhao2019pyod} provides the detector library from which
most of our detector scores are extracted.

\paragraph{Standard vision.}
\textbf{CIFAR-10}~\cite{krizhevsky2009cifar}, \textbf{ImageNet}~\cite{deng2009imagenet},
\textbf{ResNet}~\cite{he2016resnet}, \textbf{ViT}~\cite{dosovitskiy2021vit},
and \textbf{Imagenette}~\cite{howard2019imagenette} are the standard
building blocks of our online TTA and CIFAR-10-C experiments.

\paragraph{Label-free performance prediction and accuracy estimation.}
\textbf{Deng and Zheng}~\cite{deng2021labels} studied the problem of evaluating
classifier accuracy without labels, developing methods based on the agreement
between classifiers and the structure of the confidence distributions.
This work is closely related to ATC \cite{garg2022atc} and Agreement-on-the-Line
\cite{baek2022agreement} in its goals and is part of the same literature that
K-Bound departs from in its relational target (sign of difference) and
pre-adaptation certificate.

\paragraph{Anomaly/outlier detection methods.}
The anomaly detection backbone of the experiments uses standard methods from
the literature: \textbf{Isolation Forest}~\cite{liu2008isolation},
\textbf{Local Outlier Factor}~\cite{breunig2000lof},
\textbf{One-Class SVM}~\cite{scholkopf2001estimating},
\textbf{ECOD}~\cite{li2022ecod}, and the PyOD suite~\cite{zhao2019pyod}.

\section{Positioning: The Gap K-Bound Fills}\label{sec:rel-gap}

Each line of work supplies one ingredient and lacks the rest; the gap is the
combination none of them holds at once. Table~\ref{tab:rel} makes this concrete,
reproducing and expanding the positioning table from the working paper. Its
columns are the five properties jointly required for K-Bound's contribution:
label-free operation, the distinction between adapting and merely wrapping a
candidate, benefit prediction \emph{before} acting, abstention on the adaptation
decision itself, and a formal false-adapt guarantee. No prior row carries all
five checks.

\begin{table}[h]
\centering
\small
\begin{tabular}{lccccc}
\toprule
 & label- & adapts? & predicts benefit & abstains on & false-adapt \\
Work & free? & & \emph{before} adapting? & the decision? & guarantee? \\
\midrule
Tent \cite{wang2021tent} & yes & yes & no & no & no \\
EATA \cite{niu2022eata} & yes & yes & partial (filtering) & no & no \\
SAR \cite{niu2023sar} & yes & yes & stability only & no & no \\
CoTTA \cite{wang2022cotta} & yes & yes & no & no & no \\
TTT \cite{sun2020ttt} & yes & yes & no & no & no \\
MEMO \cite{zhang2022memo} & yes & yes & no & no & no \\
LAME \cite{boudiaf2022lame} & yes & yes & no & no & no \\
NOTE \cite{gong2022note} & yes & yes & no & no & no \\
ATC \cite{garg2022atc} & yes & monitors & post-hoc estimate & --- & no \\
Agr.-on-Line \cite{baek2022agreement} & yes & monitors & post-hoc estimate & --- & no \\
BBSE \cite{lipton2018detecting} & partial & corrects & corrects shift & no & no \\
\textbf{\KGA\ (ours)} & yes & wrapper & \textbf{yes} & \textbf{yes} & \textbf{yes} \\
\bottomrule
\end{tabular}
\caption{Positioning of K-Bound against adjacent methods. Prior TTA methods
adapt (some with internal guards); performance estimators operate post-hoc.
No prior method provides a single pre-adaptation adapt/freeze/abstain
certificate with an explicit, proven unknowable region and a false-adapt bound.
\KGA\ wraps any TTA candidate without modifying it.}
\label{tab:rel}
\end{table}

The surviving gap is precise: no prior method provides a single
adapt/freeze/abstain certificate with an \emph{explicit, proven} unknowable
region and a false-adapt guarantee.

To state this more carefully:
\begin{itemize}
\item TTA mechanisms adapt and (some) detect failure, but commitment happens
    before detection and no false-adapt rate is bounded.
\item Performance estimators (ATC, Agreement-on-the-Line) estimate $R_T(f)$
    post-hoc; they do not certify $\sign\Delta$ before the adaptation decision.
\item Guarded adaptation (EATA filtering, SAR resets) protects one failure
    direction but does not pose the three-way decision with an unknowable
    region.
\item Identifiability theory (Ben-David et al., BBSE) shows when risk is
    identifiable but does not yield a finite-sample, false-adapt-controlled
    operational certificate.
\item Selective prediction (Geifman and El-Yaniv) abstains on individual
    predictions, not on the adaptation decision as a whole.
\item Conformal prediction (Vovk et al., Angelopoulos and Bates) provides the
    interval mechanism K-Bound uses, but not the framing of the
    adapt/freeze/abstain decision.
\end{itemize}

K-Bound combines the identifiability insight (when is $\sign\Delta$
recoverable?) with the conformal certificate mechanism (how do we control
the false-adapt rate given a valid interval?) and the selective-abstention
logic (when the interval is too wide to certify either sign, the correct
action is to abstain) into a single coherent decision framework.
The result is a method that is simultaneously label-free, mechanism-agnostic,
pre-adaptive, three-way, and formally guaranteed --- a combination no prior
work achieves.
Closing the gap demands precise objects: a benefit whose sign is the decision,
an evidence map that sees only covariates, and a trichotomy of regimes that
classifies what such evidence can and cannot certify. The next chapter fixes
exactly these, and with them the impossibility theorem and the certificate that
the rest of the manuscript proves.
